\subsection{The realized dividend-income tilt at matched breadth}

A matched-breadth comparison isolates the dividend pattern from the
Top 10/Top 20 construction choice. Across the six original-anchor cohorts,
mean gross paid distributions as a fraction of initial capital are:

\begin{minipage}{\linewidth}
\begin{center}\small\setlength{\tabcolsep}{4pt}\renewcommand{\arraystretch}{1.18}
\begin{longtable}{@{}lrrrr@{}}\toprule
Breadth & G10 paid yield & B02 paid yield & G10 paid + unpaid & B02 paid + unpaid\\\midrule\endhead
Top 10 & 0.56\% & 2.02\% & 0.60\% & 2.22\%\\
Top 20 & 0.64\% & 1.75\% & 0.67\% & 1.91\%\\
\bottomrule\end{longtable}\end{center}
\end{minipage}\par

At matched Top 10 breadth, B02 collects 2.02\% in paid distributions
versus G10's 0.56\%---about 3.60 times as much---without an additional
dividend factor in Equation 6. Including earned unpaid distributions
raises those figures to 2.22\% and 0.60\%. This observed income tilt
is consistent with selecting more dividend-paying companies through
volatility normalization. These are cumulative holding-period cash yields
on initial capital, averaged equally across the six cohorts.
